% =====================================================================
%  tikzlib · ciencias/moeller — diagrama de Möller (regla de las diagonales)
%  Requiere: \usepackage{tikz} ; \usetikzlibrary{calc, arrows.meta}
%  Uso:   % =====================================================================
%  tikzlib · ciencias/moeller — diagrama de Möller (regla de las diagonales)
%  Requiere: \usepackage{tikz} ; \usetikzlibrary{calc, arrows.meta}
%  Uso:   % =====================================================================
%  tikzlib · ciencias/moeller — diagrama de Möller (regla de las diagonales)
%  Requiere: \usepackage{tikz} ; \usetikzlibrary{calc, arrows.meta}
%  Uso:   % =====================================================================
%  tikzlib · ciencias/moeller — diagrama de Möller (regla de las diagonales)
%  Requiere: \usepackage{tikz} ; \usetikzlibrary{calc, arrows.meta}
%  Uso:   \input{tikzlib/ciencias/moeller.tex}   y luego  \moeller
%  Preview: tikzlib/previews/moeller.png
% =====================================================================
\newcommand{\moeller}{%
\begin{tikzpicture}[x=1.05cm, y=0.82cm, font=\sffamily\small]
  % flechas diagonales (detrás)
  \foreach \sx/\sy/\ex/\ey in {1/2/0/3, 1/3/0/4, 2/3/0/5, 2/4/0/6, 3/4/0/7, 3/5/1/7}{
    \draw[-{Stealth[length=5pt]},gray,line width=1.2pt]
      ($(\sx,-\sy)+(0.30,0.30)$) -- ($(\ex,-\ey)+(-0.34,-0.34)$);}
  % nodos
  \foreach \lbl/\x/\y in {%
    1s/0/1,
    2s/0/2, 2p/1/2,
    3s/0/3, 3p/1/3, 3d/2/3,
    4s/0/4, 4p/1/4, 4d/2/4, 4f/3/4,
    5s/0/5, 5p/1/5, 5d/2/5, 5f/3/5,
    6s/0/6, 6p/1/6, 6d/2/6,
    7s/0/7, 7p/1/7}{
    \node[draw=black, fill=white, rounded corners=1.5pt, inner sep=1.6pt,
          minimum width=0.66cm, font=\sffamily\small] at (\x,-\y) {\lbl};}
\end{tikzpicture}}
   y luego  \moeller
%  Preview: tikzlib/previews/moeller.png
% =====================================================================
\newcommand{\moeller}{%
\begin{tikzpicture}[x=1.05cm, y=0.82cm, font=\sffamily\small]
  % flechas diagonales (detrás)
  \foreach \sx/\sy/\ex/\ey in {1/2/0/3, 1/3/0/4, 2/3/0/5, 2/4/0/6, 3/4/0/7, 3/5/1/7}{
    \draw[-{Stealth[length=5pt]},gray,line width=1.2pt]
      ($(\sx,-\sy)+(0.30,0.30)$) -- ($(\ex,-\ey)+(-0.34,-0.34)$);}
  % nodos
  \foreach \lbl/\x/\y in {%
    1s/0/1,
    2s/0/2, 2p/1/2,
    3s/0/3, 3p/1/3, 3d/2/3,
    4s/0/4, 4p/1/4, 4d/2/4, 4f/3/4,
    5s/0/5, 5p/1/5, 5d/2/5, 5f/3/5,
    6s/0/6, 6p/1/6, 6d/2/6,
    7s/0/7, 7p/1/7}{
    \node[draw=black, fill=white, rounded corners=1.5pt, inner sep=1.6pt,
          minimum width=0.66cm, font=\sffamily\small] at (\x,-\y) {\lbl};}
\end{tikzpicture}}
   y luego  \moeller
%  Preview: tikzlib/previews/moeller.png
% =====================================================================
\newcommand{\moeller}{%
\begin{tikzpicture}[x=1.05cm, y=0.82cm, font=\sffamily\small]
  % flechas diagonales (detrás)
  \foreach \sx/\sy/\ex/\ey in {1/2/0/3, 1/3/0/4, 2/3/0/5, 2/4/0/6, 3/4/0/7, 3/5/1/7}{
    \draw[-{Stealth[length=5pt]},gray,line width=1.2pt]
      ($(\sx,-\sy)+(0.30,0.30)$) -- ($(\ex,-\ey)+(-0.34,-0.34)$);}
  % nodos
  \foreach \lbl/\x/\y in {%
    1s/0/1,
    2s/0/2, 2p/1/2,
    3s/0/3, 3p/1/3, 3d/2/3,
    4s/0/4, 4p/1/4, 4d/2/4, 4f/3/4,
    5s/0/5, 5p/1/5, 5d/2/5, 5f/3/5,
    6s/0/6, 6p/1/6, 6d/2/6,
    7s/0/7, 7p/1/7}{
    \node[draw=black, fill=white, rounded corners=1.5pt, inner sep=1.6pt,
          minimum width=0.66cm, font=\sffamily\small] at (\x,-\y) {\lbl};}
\end{tikzpicture}}
   y luego  \moeller
%  Preview: tikzlib/previews/moeller.png
% =====================================================================
\newcommand{\moeller}{%
\begin{tikzpicture}[x=1.05cm, y=0.82cm, font=\sffamily\small]
  % flechas diagonales (detrás)
  \foreach \sx/\sy/\ex/\ey in {1/2/0/3, 1/3/0/4, 2/3/0/5, 2/4/0/6, 3/4/0/7, 3/5/1/7}{
    \draw[-{Stealth[length=5pt]},gray,line width=1.2pt]
      ($(\sx,-\sy)+(0.30,0.30)$) -- ($(\ex,-\ey)+(-0.34,-0.34)$);}
  % nodos
  \foreach \lbl/\x/\y in {%
    1s/0/1,
    2s/0/2, 2p/1/2,
    3s/0/3, 3p/1/3, 3d/2/3,
    4s/0/4, 4p/1/4, 4d/2/4, 4f/3/4,
    5s/0/5, 5p/1/5, 5d/2/5, 5f/3/5,
    6s/0/6, 6p/1/6, 6d/2/6,
    7s/0/7, 7p/1/7}{
    \node[draw=black, fill=white, rounded corners=1.5pt, inner sep=1.6pt,
          minimum width=0.66cm, font=\sffamily\small] at (\x,-\y) {\lbl};}
\end{tikzpicture}}
